% p3_work_impact.tex
% AI Agent 对未来工作可能带来的冲击
% 编译方式: xelatex p3_work_impact.tex

\documentclass[12pt,a4paper]{article}

% ============================================================
% 字体与中文支持
% ============================================================
\usepackage{fontspec}
\usepackage{xeCJK}
\setCJKmainfont[BoldFont=PingFang SC Semibold]{PingFang SC}
\setCJKsansfont[BoldFont=PingFang SC Semibold]{PingFang SC}
\setCJKmonofont{PingFang SC}
\setmainfont{Palatino}
\setsansfont{Helvetica Neue}
\setmonofont{Menlo}

% ============================================================
% 页面布局
% ============================================================
\usepackage[
  top=2.5cm,
  bottom=2.5cm,
  left=2.5cm,
  right=2.5cm
]{geometry}

% ============================================================
% 常用宏包
% ============================================================
\usepackage{amsmath,amssymb}
\usepackage{graphicx}
\graphicspath{{images/}}
\usepackage{longtable,booktabs,array}
\usepackage{enumitem}
\usepackage{xcolor}
\usepackage{hyperref}
\hypersetup{
  colorlinks=true,
  linkcolor=blue!70!black,
  urlcolor=blue!70!black,
  pdfauthor={整理自李宏毅老师 AI Agent 系列课程},
  pdftitle={AI Agent 对未来工作可能带来的冲击},
}

% ============================================================
% 段落与间距
% ============================================================
\usepackage{parskip}
\setlength{\parindent}{0pt}
\setlength{\parskip}{6pt plus 2pt minus 1pt}

% ============================================================
% 引用环境美化
% ============================================================
\usepackage{mdframed}
\newmdenv[
  topline=false,
  bottomline=false,
  rightline=false,
  linewidth=3pt,
  linecolor=blue!40,
  innerleftmargin=12pt,
  innerrightmargin=12pt,
  innertopmargin=8pt,
  innerbottommargin=8pt,
  backgroundcolor=blue!3,
  skipabove=10pt,
  skipbelow=10pt,
]{myquote}

% ============================================================
% 节标题样式
% ============================================================
\usepackage{titlesec}
\titleformat{\section}{\Large\bfseries\sffamily}{}{0em}{}[\vspace{-0.5em}\rule{\textwidth}{0.5pt}]
\titleformat{\subsection}{\large\bfseries\sffamily}{}{0em}{}
\titleformat{\subsubsection}{\normalsize\bfseries\sffamily}{}{0em}{}
\setcounter{secnumdepth}{0}  % 不编号

% ============================================================
% 防止 overfull
% ============================================================
\setlength{\emergencystretch}{3em}
\tolerance=1000

% ============================================================
% 图片插入命令
% ============================================================
\newcommand{\keyslide}[2]{%
  \begin{figure}[htbp]
    \centering
    \includegraphics[width=0.92\textwidth]{#1}
    \caption{#2}
  \end{figure}
}

% ============================================================
\begin{document}

% ============================================================
% 标题页
% ============================================================
\begin{titlepage}
  \centering
  \vspace*{3cm}
  {\Huge\bfseries\sffamily AI Agent 对未来工作可能带来的冲击\par}
  \vspace{1.5cm}
  {\Large 整理自李宏毅老师 AI Agent 系列课程第 3 集\par}
  \vspace{2cm}
  {\large\sffamily 论文撰写 \textbullet{} 模型训练 \textbullet{} 学术评估\par}
  \vfill
  {\large \today\par}
\end{titlepage}

\newpage
\tableofcontents
\newpage

% ============================================================
\section{课程概览}
% ============================================================

本集课程探讨 AI Agent 对未来工作可能带来的冲击，并以学术研究领域为例进行深入分析：

\begin{enumerate}[leftmargin=2em]
  \item \textbf{AI 独立撰写研究论文的能力与成本优势}
  \item \textbf{AI 自主训练模型与探索}
  \item \textbf{AI 产生研究想法的新颖性对比}
  \item \textbf{AI 担任论文审查委员的现状与争议}
  \item \textbf{AI 举办并参与学术会议的实验}
\end{enumerate}

% ============================================================
\section{一、AI 独立撰写研究论文的能力}
% ============================================================

\keyslide{../key_slides/p3_slide_0005.jpg}{AI写论文：The 100x Research Institution}

\subsection{从工具到自主 Agent 的演进}

AI 扮演的角色正在发生改变：
\begin{itemize}[leftmargin=2em]
  \item \textbf{最早阶段}：作为工具，一个口令一个动作。
  \item \textbf{协作阶段}：人们开始与 AI 协作，共同完成任务。
  \item \textbf{当前阶段}：AI Agent 具备了更强的自主性，有机会独立完成复杂的任务。
\end{itemize}

对于学术研究而言，最核心的问题是：\textbf{AI 能不能自己写一篇文章？}

\subsection{斯坦福教授的实验："100倍的 Research Assistant"}

斯坦福大学政治经济学教授 Andrew Hall 进行了一项实验：
\begin{itemize}[leftmargin=2em]
  \item 他让 Claude 扩充他过去针对美国大选的研究，使用新的数据和旧的分析方法重新跑一遍。
  \item Prompt 设计得非常细致，就像指导教授在教研究生做研究。
  \item \textbf{AI 的表现}：Claude 花了 1 个小时，花费约 10 美金，完成了一篇新论文。
  \item \textbf{人类对照组}：一名博士生接到同样的指令，花了 16 个小时（两个工作日），按照美国行情至少需要 1000 美金。
\end{itemize}

\textbf{结果对比}：
\begin{itemize}[leftmargin=2em]
  \item 人类做的稍微好一点点，Claude 在处理数据时犯了一个错。
  \item 但是从成本来看，AI 比人类便宜 100 倍。即使让 AI 重复运行 5 次来寻找错误，也只花费 50 美金，依然比人类便宜 20 倍。
\end{itemize}

这引发了一个思考：未来最有生产力的研究机构，可能是一个资深老师带着一群 AI Agent（"龙虾"），而不是一群人类研究生。

\subsection{AI 代劳研究的意义探讨}

有人可能会反感：研究本来就应该是人类做的事情，怎么能由 AI 代劳？

但我们需要回归研究的核心价值：\textbf{找出问题、解决问题，让世界变得更好。}
如果 AI 能够比人类更好地找出并解决问题，那么让 AI 来做研究并没有什么不对。

% ============================================================
\section{二、台湾使用行为分析与附录里的"正文"}
% ============================================================

有一篇关于台湾人使用 Claude 行为的\href{https://arxiv.org/abs/2602.17221}{分析文章}。有趣的是，这篇看似是主体的分析文章，实际上只是真正论文的\textbf{附录}。

这篇文章的"正文"，其实是在展示\textbf{如何给 Claude 写 Prompt，让它能够近乎全自动地写出一篇分析文章}。人类在其中的角色仅仅是检查。这进一步证明了 AI Agent 在处理文献收集、数据分析和论文撰写方面的高度自动化潜力。

% ============================================================
\section{三、AI 自主训练模型与实验}
% ============================================================

除了文献和数据分析，在需要建模型、跑实验的领域，AI Agent 也能胜任吗？

Andrej Karpathy 释出的 \textbf{Auto Research} 展示了惊人的能力：
\begin{itemize}[leftmargin=2em]
  \item 让一个 LLM 自动帮你训练模型。
  \item Agent 大约每五分钟进行一次实验（"心跳一次"），不断调整 training script。
  \item 在整个过程中\textbf{完全没有人类介入}。
  \item AI 先训练第一版模型，观察结果，思考需要修改的地方，然后训练第二版、第三版……模型表现越来越好。
\end{itemize}

这证明了 AI 具有自主进行实验迭代和优化的能力。

% ============================================================
\section{四、AI 产生研究想法的能力}
% ============================================================

\keyslide{../key_slides/p3_slide_0010.jpg}{AI写论文：产生研究想法}

也许 AI 能写文章、做实验，但"寻找问题"、"想出新颖的 Idea"总该是人类的专利吧？

\subsection{2024年的研究：AI 点子更具创新性？}

\href{https://arxiv.org/abs/2409.04109}{一篇 2024 年的论文}探讨了 LLM 能否产生新颖的研究 Idea：
\begin{itemize}[leftmargin=2em]
  \item 让语言模型对过去的论文进行 RAG（检索增强生成），产生大量研究想法。
  \item 找人类专家也提出研究想法。
  \item 找另一群专家进行盲测评价（包括新颖性、可行性、有效性等指标）。
\end{itemize}

\textbf{结果出人意料}：在多数指标上，AI 击败了人类。人类唯一赢过 AI 的是\textbf{可行性（Feasibility）}，但在\textbf{新颖性（Novelty）}上，专家认为 AI 想出的题目更具创新。

\subsection{2025年的续作：实作后的真相}

一年后，同一个团队进行了续作研究：将 AI 和人类提出的点子真正做成短篇论文，再次进行评价。

\textbf{结果反转}：
\begin{itemize}[leftmargin=2em]
  \item 当 AI 的点子真正被实作后，其新颖性评分大幅下降，最终比不上人类。
  \item 原因在于：AI 有时会"堆砌新颖词汇"，表面上看起来很厉害，但实际执行时才发现做不起来。
  \item \textbf{结论}：至少在当时的模型能力下，产生好问题的能力，人类依然占据优势。
\end{itemize}

% ============================================================
\section{五、AI 担任论文审查委员（Reviewer）}
% ============================================================

不仅是写论文，AI 也开始参与论文审查流程。

在某个与 AI 相关的国际会议（如 AAAI）中，AI 正式进入了审查流程：
\begin{itemize}[leftmargin=2em]
  \item 每篇文章不仅有人类 Reviewer，还会明确配备一个 AI Reviewer（它会标明自己是 AI）。
  \item AI 不打分数，只给意见供人类参考。
  \item 此外，还有 AI 担任 Meta Reviewer。
\end{itemize}

讲者分享了自己担任 Area Chair 时的经历：他发现有些标明为人类的 Reviewer，其评审意见的第一句话竟然是 "Sure, I can help you write this review"，这表明许多"人类 Reviewer"背后其实也是 AI Agent。

\subsection{使用 AI 审查的原则}

讲者并不反对使用 AI 辅助审查，因为 Review 的核心意义是\textbf{找出文章的问题并使其更好}。如果 AI 能比人类更敏锐地发现问题，那由 AI 审查对作者更有帮助。

但他反对的是\textbf{使用不够好的 AI 模型}进行审查。如果 AI 给出"牛头不对马嘴"的评审意见，并且在被指出错误后只是敷衍修改，那是无法接受的。

讲者提到自己指导的 AI Agent "小金" 也经常帮学生审查论文。通过精心设计的 Prompt，小金的审查不仅能指出问题，还能给出具体建议；并且会根据截止日期的远近，调整审查的严格程度，甚至在临近 Deadline 时提供情绪价值的鼓励。

% ============================================================
\section{六、AI Agent for Science：全 AI 学术会议实验}
% ============================================================

\keyslide{../key_slides/p3_slide_0015.jpg}{AI Agent for Science 会议}

既然 AI 可以写论文，也可以审查论文，这就形成了一个闭环。

斯坦福的研究人员举办了一个名为 \textbf{AI Agent for Science} 的实验性会议：
\begin{itemize}[leftmargin=2em]
  \item \textbf{规则}：论文的主要贡献者和第一作者必须是 AI，并由 AI 审查论文。
  \item 投稿时必须标明人类的介入程度，要求人类介入越少越好。
  \item \textbf{结果}：接收率极低（247 篇中仅接收 48 篇，小于 20\%）。
\end{itemize}

\textbf{核心发现}：
\begin{itemize}[leftmargin=2em]
  \item 那些最终被接收的优秀论文，其\textbf{点子发想}和\textbf{实验设计}阶段，人类的介入程度明显较高。
  \item 而在\textbf{数据分析}和\textbf{论文写作}阶段，AI 几乎可以独立完成。
\end{itemize}

投稿者的反馈也印证了这一点：目前的 AI 仍然难以想出真正有创造力的新颖点子，多数时候只是对已有事物的重组。

% ============================================================
\section{总结}
% ============================================================

\begin{itemize}[leftmargin=2em]
  \item \textbf{AI 的能力边界}：AI 已经具备了收集文献、分析数据、撰写论文甚至自主运行实验的能力，并且成本极低。
  \item \textbf{人类的核心价值}：目前 AI 在提出真正新颖且可行的研究问题上仍显不足。人类在发想好问题、引导研究方向上的角色依然不可替代。
  \item \textbf{人机协作新模式}：未来的研究模式可能是人类负责"点子发想"和"实验设计"（即"出主意"），而由 AI Agent（"龙虾"们）去完成数据分析和论文写作等繁重工作。
\end{itemize}

\end{document}
